%% MoCHII, M3, and MAPPINGS V 5.2.1: algorithmic and physical comparison.
%% Author: Kwang-Il Seon
%% Written for my_memo.cls.
\documentclass[english,a4paper]{my_memo}
\synctex=1
\usepackage{amsmath}
\usepackage{amssymb}
\usepackage{microtype}
\usepackage{booktabs}
\usepackage{longtable}
\usepackage{array}
\usepackage{enumitem}
\hypersetup{hidelinks}

\newcommand{\code}[1]{\texttt{#1}}
\newcommand{\MoCHII}{\textsc{MoCHII}}
\newcommand{\Mthree}{M$^3$}
\newcommand{\Mappings}{MAPPINGS V 5.2.1}
\newcommand{\HII}{H\,\textsc{ii}}
\newcommand{\HI}{H\,\textsc{i}}
\newcommand{\HeI}{He\,\textsc{i}}
\newcommand{\HeII}{He\,\textsc{ii}}
\newcommand{\HeIII}{He\,\textsc{iii}}
\newcommand{\Te}{T_{\rm e}}
\newcommand{\nelec}{n_{\rm e}}
\newcolumntype{P}[1]{>{\raggedright\arraybackslash}p{#1}}

\begin{document}

\title{MoCHII, M$^3$, and MAPPINGS V 5.2.1:\\
An Algorithmic and Physical Comparison}
\author{Kwang-Il Seon}
\date{\docmoddate}
\maketitle

\begin{abstract}
This memo compares \MoCHII{} (MOnte Carlo for \HII{} regions), the public
\Mthree{} (Messenger Monte Carlo MAPPINGS V) tree at
\code{\char`~/RT\_Codes/MAPPINGS/M3-main}, and the independent public
\Mappings{} release at
\code{\char`~/RT\_Codes/MAPPINGS/mappings\_V-521}.  The comparison is
based on connected execution paths and quantitative tests rather than on
routine names alone.  \MoCHII{} is strongest as a modern
three-dimensional transport and synthetic-observation engine: octree AMR,
an analytic absorption estimator, a dedicated Cartesian DDA backend,
solution-driven ionization-front refinement, dust absorption and
scattering, stochastic dust/PAH emission, and observer peel-off imaging.
\Mappings{} is strongest as a deterministic one-dimensional general
plasma engine: a 30-element option, density-dependent multilevel cooling,
finite-source-age photoionization, non-equilibrium cooling, radiation
pressure, and steady radiative MHD shocks with an iterated precursor.
Its 1D photoionization path also connects detailed grain/PAH physics to a
stochastic infrared solver.  \Mthree{} occupies a distinct middle ground:
it grafts three-dimensional Monte Carlo transport onto an older MAPPINGS
V 5.1.13-era physics tree and therefore combines broad MAPPINGS
microphysics with arbitrary Cartesian geometry, but it does not inherit
all later 5.2.1 changes or expose all parent modes through its public 3D
driver.  The principal limitations are correspondingly different.
\MoCHII{} presently uses low-density-limit metal cooling fits in its hot
thermal loop and has an \HII-region-centered ion inventory; \Mthree{} has
no AMR and a communication-heavy uniform-grid packet algorithm; and
\Mappings{} is restricted to spherical or plane-parallel 1D structures,
uses an outward-only diffuse-field approximation, and retains a legacy
interactive global-state architecture.  The recommended strategy is to
retain the \MoCHII{} transport architecture, use \Mappings{} as the
current microphysics and 1D shock/time-dependence reference, and use
\Mthree{} specifically to study how MAPPINGS physics behaves after being
mapped into three dimensions.
\end{abstract}

\tableofcontents

%==============================================================================
\section{Purpose, scope, and basis of the comparison}

\MoCHII{} and \Mthree{} are self-consistent three-dimensional Monte Carlo
photoionization codes, whereas \Mappings{} is a deterministic one-dimensional
plasma and radiative-shock code.  \MoCHII{} adds gas
physics to the validated MoCafe v2.00 dust radiative-transfer engine.  Its
founding design treats grid traversal, radiation estimators, dust physics,
parallelism, file formats, and synthetic observations as first-class
components.  \Mthree{} adds three-dimensional Monte Carlo radiation
transport to MAPPINGS V, whose principal strength is its extensive atomic,
ionic, heating, cooling, and emission-line physics \cite{jin2022,sd1993}.
\Mappings{} continues that parent line as version 5.2.1, with separate
photoionization, shock, equilibrium-cooling, non-equilibrium-cooling, and
single-ion diagnostic drivers.

This memo distinguishes three levels of evidence:

\begin{enumerate}[leftmargin=2em]
\item \emph{Physics present in the parent code.}  The MAPPINGS V source
  tree contains one-dimensional photoionization, shock, non-equilibrium,
  dust, and plasma utilities.
\item \emph{Physics connected to the three-dimensional driver.}  A routine
  appearing in the source tree is not automatically active in the
  \code{montph7} Monte Carlo path.
\item \emph{Physics quantitatively gated in the current code.}  This is the
  strongest evidence: an analytic test, independent solver, database
  comparison, or published benchmark with a stated numerical tolerance.
\end{enumerate}

This distinction is essential for shocks and time dependence.  The public
MAPPINGS tree contains \code{shock2--shock5}, \code{timion}, and
\code{timtqui}; however, the three-dimensional \code{montph7} family sets
\code{nmod='EQUI'} before the cell thermal solve.  Thus the present
three-dimensional \Mthree{} calculation is an equilibrium photoionization
model.  The legacy routines remain valuable MAPPINGS capabilities, but
they should not be described as coupled three-dimensional \Mthree{}
physics without additional development and validation.

No direct same-input, same-atomic-data timing or line-flux campaign among
the three local executables has yet been performed.  Numerical performance
comparisons in this memo therefore use algorithm structure and each code's
own recorded gates.  They are diagnostic rather than a substitute for a
controlled cross-code experiment.

%==============================================================================
\section{Lineage: why MAPPINGS V 5.2.1 is not
\texorpdfstring{M$^3$}{M3} 5.2.1}

The three code names can easily obscure the most important architectural
fact.  \Mthree{} is not a three-dimensional front end to the newly added
\Mappings{} directory.  Its source identifies itself as MAPPINGS V
5.1.13 with a separate Monte Carlo development history.  The two trees
still share 78 Fortran filenames, but the shared implementations have
diverged substantially; the 5.2.1 versions of core files such as
\code{photo6.f}, \code{photo7.f}, \code{shock5.f}, \code{mapinit.f}, and
\code{cool.f} contain extensive later changes.  \Mthree{} adds its own
\code{montph7}, Cartesian packet transport, PDF, and three-dimensional
thermal-solver files instead \cite{jin2022,mappings521}.

This gives each branch a different meaning:

\begin{itemize}[leftmargin=2em]
\item \Mappings{} is the newer parent-physics branch.  Its public
  \code{map52} executable offers P6/P7 multizone photoionization, S5
  shock plus precursor models, CIE and NEQ cooling curves, PIE curves,
  single slabs, and multilevel-ion diagnostic modes.
\item \Mthree{} is the three-dimensional geometry branch.  It is the
  relevant comparison for arbitrary density fields, but its physics
  inventory must be assessed from the connected \code{montph7} path,
  not inferred from the presence of a similarly named MAPPINGS routine.
\item \MoCHII{} is an independent three-dimensional transport branch built
  from MoCafe rather than MAPPINGS.  Similar physical terms therefore have
  independent algorithms, atomic-data pipelines, and validation histories.
\end{itemize}

Consequently, a discrepancy between \MoCHII{} and \Mthree{} cannot safely
be labeled a discrepancy with \Mappings{} 5.2.1.  For future differential
tests, 5.2.1 should be the primary one-zone and 1D microphysics reference;
\Mthree{} should be treated as the older MAPPINGS-derived 3D realization.

%==============================================================================
\section{MAPPINGS V 5.2.1 in its own right}

\subsection{Deterministic 1D integration and diffuse radiation}

P6 and P7 integrate successive zones in spherical or plane-parallel
geometry.  Step sizes respond to optical depth and predicted changes in
temperature and ionization.  The result has no Monte Carlo shot noise and
is therefore efficient for large parameter grids, detailed spectra, and
thin cooling zones.  Isobaric models can include radiation pressure,
$P(r)=P_0+P_{\rm rad}(r)$, while other choices provide constant density or
a prescribed density law.  The model may terminate at a distance, an ion
or hydrogen column, or an optical/radiation-bounded condition.

The price is geometric closure.  The diffuse field is explicitly labeled
``OUTWARD ONLY'' in \code{photo6.f} and \code{photo7.f}.  It follows the
allowed 1D directions rather than transporting recombination photons over
arbitrary solid angle.  This is fast and often adequate for a monotonic
sphere or slab, but it cannot represent porous escape, oblique shadows,
clump illumination, or inward/backward diffuse coupling as a true 3D
transport calculation can.  The word ``adaptive'' in the S5 menu refers
to an adaptive one-dimensional shock integration, not to a 3D AMR mesh.

\subsection{Time dependence, cooling flows, and shocks}
\label{sec:mapstrength}

Unlike the public \Mthree{} 3D driver, the P6/P7 path genuinely connects a
finite source lifetime to the time-dependent ion solver.  It supports an
equilibrium solution, a finite-age ionization calculation, and a
post-equilibrium source-off evolution.  The implementation advances ionic
populations with adaptive substeps and couples the elapsed ionization-front
time to each new zone.  This is a useful one-dimensional ionization-history
model, although it is not radiation hydrodynamics and does not evolve an
arbitrary multidimensional density field.

The S5 path is an even clearer 5.2.1 advantage.  It solves a steady,
plane-parallel radiative MHD shock, adapts its step to atomic, cooling, and
photon-absorption timescales, and iterates the shock radiation with an
automatically calculated photoionized precursor.  The source description
advertises shock speeds from roughly 10 to $10^4$ km s$^{-1}$.  Companion
drivers compute CIE cooling, non-equilibrium cooling, and PIE curves.  For
one-dimensional shock spectra and cooling flows, neither current 3D code
offers an equivalent connected and mature path.

\subsection{Atomic, electron, and dust physics}

The standard preference file activates 16 elements, while
\code{mapFull.prefs} activates all 30 elements from H through Zn; static
dimensions permit 31 ion stages and 10,240 photon-energy entries.  The
thermal loop directly calls fine-structure, intercombination, resonance,
multilevel, iron, recombination, free-free, collisional-ionization,
charge-exchange, grain/PAH, and cosmic-ray terms.  Optional $\kappa$
electron distributions modify excitation, ionization, and recombination
rates.  These are major strengths for broad plasma diagnostics and for
gas near or above coolant critical densities.

The 1D dust path is materially more complete than the public \Mthree{} 3D
dust branch.  In addition to grain charge, depletion, photoelectric
heating, gas-grain collisions, PAHs, and destruction prescriptions,
P6/P7 call \code{dusttemp} zone by zone.  That routine implements a
Draine--Li stochastic temperature distribution and returns an infrared
photon field.  Thus \Mappings{} has a coupled 1D gas--dust--IR calculation,
whereas \MoCHII{} has the stronger spatially resolved 3D dust transport
and imaging chain.

Several cautions remain.  The atomic tables are broad but globally coupled
and are harder to provenance or replace ion by ion than the \MoCHII{}
registry.  A Compton routine exists but no call from the examined main
cooling path was found, and the microturbulent heating block in
\code{cool.f} is commented out; neither should be counted as active 5.2.1
physics without a dedicated execution test.  Resonance-line attenuation
uses escape/transfer approximations rather than a multidimensional,
frequency-redistributing Monte Carlo line solver.

\subsection{Software and workflow limitations}

\Mappings{} preserves a productive but legacy Fortran architecture:
common blocks, compile-time maxima, runtime table loading, interactive
menus, and many ASCII output products.  P7 exposes experimental routines
and an API, but the normal workflow remains less scriptable and less
modular than \MoCHII{}'s namelist, HDF5/FITS, module, and Python-reader
design.  The documentation overview describes the manual as work
in progress, so reliable use still requires source reading.  The benefit
of this conservatism is continuity with decades of MAPPINGS models; the
cost is a higher barrier to automated regression testing, embedding, and
isolated replacement of atomic data.

%==============================================================================
\section{Executive comparison}

\begin{longtable}{P{0.19\textwidth}P{0.35\textwidth}P{0.35\textwidth}}
\caption{High-level comparison of the current three-dimensional codes.}
\label{tab:executive}\\
\toprule
Topic & \MoCHII{} & \Mthree{} \\
\midrule
\endfirsthead
\multicolumn{3}{l}{\small\itshape Table~\ref{tab:executive} continued.}\\
\toprule
Topic & \MoCHII{} & \Mthree{} \\
\midrule
\endhead
\bottomrule
\endfoot
Grid & Octree AMR or a memory-minimal uniform Cartesian backend; optional
solution-driven I-front refinement. & Uniform Cartesian grid divided into
equal MPI blocks; no AMR. \\

Absorption transport & Analytic leaf-by-leaf attenuation estimator when
the band has no scattering; only direction and frequency are sampled. &
Lucy-style random optical depth, discrete absorption, and immediate local
re-emission. \\

Diffuse field & Explicit selected H/He recombination channels rebuilt from
the current state; case B OTS remains available. & Local continuum and line
emissivity PDF sampled immediately after gas absorption; isotropic
re-emission. \\

Frequency treatment & Dynamically allocated FUV and EUV segments, normally
tens of rate-optimized bins with threshold alignment. & Up to 10,240 bins;
a fixed number of packets is launched at every frequency and packet energy
follows the source spectrum. \\

Parallel model & Packet decomposition and shared-memory grid arrays; one
full grid copy per node. & Spatial domain decomposition; packets crossing
blocks are synchronized through MPI. \\

Gas microphysics & Modern, provenance-bearing fitted data for an
H\,\textsc{ii}-region/PDR-centered registry; low-density cooling fits in
the thermal loop. & Broad MAPPINGS ion network and multilevel line cooling
inside the thermal loop; older but extensively exercised atomic-data
compilation. \\

Dust & EUV/FUV absorption and HG scattering, live dust/PAH survival,
equilibrium temperature, SEDust stochastic emission, and IR maps. & Rich
grain charge, depletion, photoelectric, and PAH microphysics; the public 3D
dust transport branch treats dust extinction/absorption but does not expose
an equivalent scattering plus stochastic-IR imaging chain. \\

Synthetic observations & Leaf emissivities, flux-conserving line and field
maps, dust-band maps, multi-observer peel-off images, and spectral cubes. &
Extensive line spectrum and cell-resolved emissivities; maps can be made by
projection, but no equivalent general observer peel-off layer is evident in
the public path. \\

Best current use & Dusty, clumpy, spatially resolved H\,\textsc{ii}
regions and atomic PDR skins on simulation-derived density fields. & Broad
nebular line diagnostics, denser gas, harder PN-like spectra, and problems
that benefit from the full MAPPINGS ion and cooling inventory. \\
\end{longtable}

The concise verdict is:

\begin{quote}
\MoCHII{} is presently the stronger three-dimensional transport, dust, AMR,
and synthetic-observation framework.  \Mthree{} is presently the stronger
broad-spectrum nebular microphysics framework, especially when many ion
stages or density-dependent line cooling determine the solution.
\end{quote}

\begin{longtable}{P{0.24\textwidth}P{0.34\textwidth}P{0.32\textwidth}}
\caption{What \Mappings{} adds to the original two-code comparison.}
\label{tab:mappings}\\
\toprule
Topic & \Mappings{} strength & Limitation relative to the 3D codes \\
\midrule
\endfirsthead
\multicolumn{3}{l}{\small\itshape Table~\ref{tab:mappings} continued.}\\
\toprule
Topic & \Mappings{} strength & Limitation relative to the 3D codes \\
\midrule
\endhead
\bottomrule
\endfoot
Photoionization & Noise-free adaptive 1D integration; equilibrium,
finite-age, and source-off evolution; radiation-pressure options. & Only
spherical or plane-parallel structure; outward-only diffuse field. \\

Shocks and cooling flows & Connected steady radiative MHD shock plus
iterated precursor, CIE cooling, NEQ cooling, and PIE modes. & Steady or
prescribed 1D histories, not multidimensional radiation hydrodynamics. \\

Microphysics & Full 30-element preference set, many ion stages,
density-dependent line cooling, $\kappa$ electrons, cosmic rays, and broad
spectral output. & Global tables are harder to audit and update ion by ion;
not every dormant routine is connected to the main cooling path. \\

Dust & Grain/PAH charge and destruction, gas--grain exchange, stochastic
temperature distribution, and a coupled 1D IR field. & No arbitrary 3D
dust geometry, shadowing, AMR, or observer image construction. \\

Numerics and workflow & Fast deterministic parameter grids with decades
of model heritage. & Legacy interactive common-block architecture,
compile-time maxima, fragmented ASCII products, and incomplete manuals. \\

Best current use & Integrated spectra of 1D H\,\textsc{ii} regions, PNe,
AGN slabs, radiative shocks, and equilibrium/non-equilibrium cooling
sequences. & Not a replacement for 3D modeling when geometry is part of
the physical question. \\
\end{longtable}

%==============================================================================
\section{Radiative-transfer algorithms}

\subsection{Radiation-field estimators}

Both three-dimensional codes estimate the cell radiation field from packet trajectories, but
their estimators differ fundamentally.  \Mthree{} follows the Lucy path
length estimator.  At frequency $\nu$, a packet with energy
$\epsilon_\nu$ crossing a distance $l$ in cell volume $V$ contributes
schematically
\begin{equation}
 J_\nu \;=\; \frac{1}{4\pi\,\Delta t\,V}
       \sum_{\rm paths}\epsilon_\nu l .
 \label{eq:m3j}
\end{equation}
The distance to the next absorption is sampled from
\begin{equation}
 \tau_p=-\ln p,
 \label{eq:taup}
\end{equation}
with $p$ uniform on $(0,1)$.  The packet crosses Cartesian cells until its
accumulated optical depth reaches $\tau_p$, is absorbed, and is re-emitted
in situ from the local emissivity distribution.  This is conceptually
general and treats diffuse radiation naturally.

For the absorption-only ionizing band, \MoCHII{} instead integrates the
attenuated packet weight analytically through every crossed leaf:
\begin{equation}
 \Delta J_b({\rm leaf}) = w L_{\rm pkt}
   \frac{e^{-\tau_{\rm in}}-e^{-\tau_{\rm out}}}{\kappa_b} .
 \label{eq:edge}
\end{equation}
The packet is followed to the domain edge and contributes the expectation
value of the absorbed path in each leaf.  There is no sampled absorption
location; statistical noise comes only from direction and bin sampling.
The analytic attenuation gate measured a mean photo-rate bias of
$+0.004\%$.  An optional Owen-scrambled Sobol launch
(\code{launch\_sequence='sobol'}) stratifies exactly those remaining
launch variables: on the same gate the cell-wise photo-rate error then
scales as $\sim 1/N$ rather than $1/\sqrt{N}$, with measured effective
photon gains of $29\times$ at $2^{17}$ packets and $222\times$ at
$2^{20}$, and the launch set becomes independent of the MPI task count.  When dust scattering is enabled, \MoCHII{} changes to a
forced-first-interaction estimator, applies the dust albedo as a packet
weight, samples an HG scattering direction, and uses Russian roulette at
low weight.

Equation~(\ref{eq:edge}) is a substantial advantage for the present
\MoCHII{} target regime: ionizing radiation in H\,\textsc{ii} regions is
normally absorption-dominated.  \Mthree{} pays the variance of discrete
absorption events even in that limit.  Conversely, the \Mthree{} packet
cycle is more general for a diffuse spectrum containing many continuum and
line channels, whereas the current \MoCHII{} diffuse emitter is explicitly
limited to selected recombination channels.

\Mappings{} has neither Monte Carlo estimator.  Its zone integrator is
deterministic and therefore essentially noise-free, but achieves that
efficiency by imposing spherical or plane-parallel symmetry and an
outward-only closure for the diffuse field.

\subsection{Diffuse radiation and energy packets}

In \Mthree{}, gas absorption converts the stellar packet into a diffuse
packet immediately.  Its new frequency is drawn from a local emissivity
PDF containing recombination continua and lines, and its new direction is
isotropic.  This implements a full three-dimensional alternative to
outward-only diffuse-field approximations \cite{jin2022}.

\MoCHII{} offers two controlled limits.  Case B uses on-the-spot
recombination.  With the explicit diffuse field enabled, case A is used
and the emission table is rebuilt each iteration from the current state.
The principal channels are H\,\textsc{ii}, He\,\textsc{ii}, and
He\,\textsc{iii} ground recombination continua; the optional He\,\textsc{i}
excited channel adds the 19.82-eV, 584-\AA{}, and two-photon contributions.
These packets use the same transport estimator as stellar packets.  This
split makes case-A/case-B bracketing and energy accounting transparent,
but it is not yet the full local emissivity inventory sampled by \Mthree{}.

\subsection{Spectral discretization}

The public \Mthree{} compile-time ceiling is 10,240 frequency entries.
The Monte Carlo loop launches $N$ packets for every frequency interval;
packet energy is proportional to $L_\nu\Delta\nu/N$.  This stratified
sampling guarantees representation of weak spectral intervals and is
attractive for detailed continua.  Its cost scales approximately with
$N_\nu N$ rather than with one total packet budget, and its large
$J_\nu(x,y,z)$ and opacity arrays are major memory consumers.

\MoCHII{} samples a frequency bin from the luminosity CDF and assigns every
packet the same band luminosity.  Its bins are designed around rate
integrals rather than spectral oversampling.  The ionizing grid can align
edges with H, He, and active metal thresholds, while an independent FUV
segment preserves the EUV resolution.  This is efficient for ionization
and thermal balance, but a fine line-rich diffuse spectrum would require
many more bins and packets than the normal 32--64-bin production setup.

\subsection{Velocity fields and line transport}

\Mthree{} reads a three-dimensional velocity field and applies a projected
Doppler shift to the packet frequency bin during continuum transport.
\MoCHII{} currently has no velocity-dependent band transport.  This is a
real \Mthree{} advantage for bulk-flow effects on continuum opacity.

None of the three codes should yet be treated as a general resonance-line transfer
solver.  \Mthree{} includes resonance emissivities but, in the published
implementation, assigns resonance lines the same cross section as the
continuum.  This is not a frequency-redistributing resonant-scattering
calculation.  \MoCHII{} explicitly treats its diagnostic lines as optically
thin emissivities and leaves resonance transfer to a dedicated code such
as LaRT.  \Mappings{} applies geometry-dependent escape and attenuation
approximations rather than multidimensional redistribution.  For forbidden
and recombination lines this is normally adequate;
for Ly$\alpha$, C\,\textsc{iv}, Mg\,\textsc{ii}, or line profiles in moving
media, neither current path is sufficient.

%==============================================================================
\section{Grid algorithms, resolution, and parallelism}

\subsection{Octree AMR versus a uniform Cartesian mesh}

The primary spatial difference is adaptive resolution.  \Mthree{} accepts
an arbitrary density distribution on a Cartesian mesh, but every cell at a
given global resolution exists.  The published HII20, HII40, and PN150
tests used $33^3$ and $55^3$ grids.  Most integrated lines agreed within
10\%, but I-front and boundary tracers such as [O\,\textsc{i}],
[O\,\textsc{ii}], and [S\,\textsc{ii}] could differ by more than 10\%
\cite{jin2022}.  This is the expected weakness of a uniform grid: the
thin zone that most needs resolution occupies a small part of the volume.

\MoCHII{} can ingest a pre-existing octree or rebuild the octree after an
initial solution.  A leaf is refined when it lies in the I-front or sees a
large face-neighbor ionization jump.  In the recorded gate, refinement
from a level-5 start to level 7 produced 464,248 leaves instead of the
2,097,152 cells of a complete level-7 grid; the effective ionized radius
agreed with the native level-7 result by $0.008\%$.  The refinement follows
clump skins and distributed fronts in turbulent media, not merely the
source position.

For problems that truly require a uniform mesh, \MoCHII{} also has a
raster-indexed Cartesian backend.  It allocates no octree topology,
neighbor table, center array, or leaf maps.  The dedicated incremental
Amanatides--Woo DDA precomputes $t_{\max}$ and $t_\Delta$ per ray and
advances by comparisons and additions.  In the recorded Str\"omgren test
it was twice as fast as the octree traversal, while matching the solution
to rounding.  At $512^3$, omitting the tree and geometry arrays saves
approximately 10 GB before physical fields are counted.

\subsection{MPI decomposition and scaling trade-offs}

\Mthree{} spatially decomposes the Cartesian domain into equal blocks.
Each rank stores the local thermal and ionization state, and packets
crossing a block boundary are transferred through collective MPI
operations.  Spatial decomposition is attractive when no node can hold
the full grid.  It also introduces frequent synchronization: the current
transport implementation collectively reduces large packet-state arrays
containing flags, positions, directions, frequencies, and absorption
counts.  At large $N_\nu$, packet count, and MPI task count, communication
can dominate.

\MoCHII{} decomposes packets rather than space.  The octree and leaf fields
reside in MPI-3 shared memory, so ranks on one node share one copy; each
node nevertheless holds a complete grid.  Ranks transport independent
packets and reduce the accumulated radiation tally at the end of a pass.
This minimizes packet communication and is well matched to Monte Carlo
transport, but the requirement of a full grid on each node ultimately
limits the largest possible AMR problem.  Thus the scaling comparison is not
one-sided: \MoCHII{} should normally scale more cleanly in packet number,
whereas \Mthree{} can distribute spatial memory more aggressively.
\Mappings{} avoids the 3D memory problem entirely by advancing one 1D
structure, which is a decisive throughput advantage only when its symmetry
assumptions are physically defensible.

%==============================================================================
\section{Ionization, thermal balance, and atomic data}

\subsection{Element and ion inventories}

The local \Mthree{} default atom file activates 16 elements:
H, He, C, N, O, Ne, Na, Mg, Al, Si, S, Cl, Ar, Ca, Fe, and Ni.  Static
array dimensions allow up to 30 elements and 31 ion stages.  The published
description characterizes the parent MAPPINGS V system as containing 30
elements and roughly 80,000 cooling and recombination lines
\cite{jin2022}.  The distinction between the parent capacity and the
default public data set should be retained when quoting these numbers.
The new \Mappings{} audit resolves the parent-code side directly:
\code{mapStd.prefs} is the 16-element setup and \code{mapFull.prefs}
activates 30 elements from H through Zn.  This full preference set, rather
than the \Mthree{} default, is the relevant breadth reference.

\MoCHII{} currently treats H and He plus an eleven-element metal registry:
C, N, O, Ne, S, Ar, Mg, Fe, Si, Cl, and Ca.  The tracked stages are chosen
for H\,\textsc{ii} regions and atomic PDR skins, rather than extending to
the highest charge state of every nucleus.  New ions are added through
provenance-bearing element, cooling-fit, and $n$-level files.  This is much
easier to audit and extend than the global MAPPINGS tables, but the present
inventory is less suitable for hard PN, X-ray, or coronal-ion problems.

\subsection{Photoionization, recombination, and charge exchange}

\MoCHII{} uses exact hydrogenic ground-state cross sections for \HI{} and
\HeII{}, and Verner et al.\ \cite{verner1996} fits for \HeI{} and metals.
Hydrogen and helium recombination defaults to the Badnell (2023) radiative
total minus the Mao--Kaastra (2016) ground-state coefficient (case B), with
the Hui--Gnedin fits retained as an option; metal RR and DR
are derived primarily from Badnell/CHIANTI v11.0.2 data.  Charge exchange
uses the modern first-stage rates of Huang et al.\ and the higher-stage
Kingdon--Ferland compilation.  Each generated file records its source,
fit range, formula, and maximum fit error.

\Mthree{} uses the established MAPPINGS compilation.  The published code
description identifies photoionization data from Seaton, Raymond,
Osterbrock, and Gould, including Auger corrections, and separates hydrogenic from
non-hydrogenic recombination.  Much of this compilation is older than the
\MoCHII{} sources, but it spans more ions and has been exercised in many
MAPPINGS applications.  ``Newer'' is not automatically equivalent to
``more accurate'': line predictions can differ by several to tens of
percent among reputable collision-strength and $A$-value compilations,
especially for iron.  The decisive test is therefore ion-by-ion
comparison with direct database evaluation and observation, not vintage
alone.

\subsection{Electron distributions and secondary processes}

Both \Mthree{} and \Mappings{} support either a Maxwellian or a $\kappa$
electron energy distribution, with rate corrections supplied by KAPPA
data.  This is
useful when exploring non-Maxwellian interpretations of nebular
diagnostics.  \MoCHII{} assumes a Maxwellian distribution.

All three codes include secondary-ionization concepts.  \MoCHII{} optionally
partitions the energy of photoelectrons above 40 eV among heat and
secondary H/He ionizations following Shull--van Steenberg.  The MAPPINGS
ion-balance routines call a secondary-ionization module and also include
cosmic-ray ionization/heating.  The MAPPINGS-derived codes therefore have
the broader general-plasma framework.  A Compton routine is present in
both source trees, but no call from either examined main cooling path was
found; it should not be counted as active physics without a dedicated test.

\subsection{The critical difference: density-dependent metal cooling}
\label{sec:densitycool}

This is the most important gas-physics advantage of both MAPPINGS-derived
codes over \MoCHII{}.  Their thermal solvers call fine-structure,
intercombination, multilevel, and iron line
modules during every evaluation of the net cooling function.  Level
populations and collisional de-excitation therefore affect the equilibrium
temperature when $\nelec$ approaches or exceeds a transition's critical
density.

\MoCHII{} deliberately separates two products:

\begin{itemize}[leftmargin=2em]
\item analytic Tier-1 cooling fits, evaluated cheaply inside the thermal
  bisection in the optically thin low-density limit; and
\item Tier-2 $n$-level statistical-equilibrium data, evaluated on the
  converged state to produce density-dependent diagnostic emissivities.
\end{itemize}

The output line ratios therefore respond to density, but collisional
quenching does not feed back into the hot thermal loop.  At ordinary
H\,\textsc{ii}-region densities this is often acceptable.  In compact
H\,\textsc{ii} regions, dense PN, shocked cooling zones, or any model with
$\nelec\gtrsim n_{\rm crit}$ for a dominant coolant, \Mthree{} is
physically better founded.  Moving the generic \MoCHII{} $n$-level solver
or a reduced density-suppression fit into the thermal loop is the
highest-value gas-physics improvement suggested by this comparison.

\subsection{Thermal processes and equilibrium solution}

All three codes solve ionization and temperature under the current
radiation field, iterating the spatial solution where required.  \MoCHII{}
brackets the root of heating minus cooling in
$\log \Te$ and re-solves the H/He and metal ionization at every trial
temperature.  Its current budget includes photoionization heating,
H/He recombination and collisional losses, free-free emission with a
tabulated Gaunt integral (H/He and, by default, the trace metals with
their net ion charge), metal-line cooling, optional metal
photoheating, secondary ionization, and the atomic-PDR grain and neutral
collision terms.

The MAPPINGS thermal solver used by \Mappings{} and inherited by \Mthree{}
includes a broader traditional plasma budget:
hydrogen and helium line cooling, fine/intercombination/resonance and
multilevel metal cooling, free-free, collisional ionization,
recombination, photoionization heating, charge-exchange energy exchange,
cosmic rays, grain/PAH photoelectric terms, and grain collisional cooling.
The microturbulent dissipation block in the examined 5.2.1
\code{cool.f} is commented out and is therefore not listed as active.
This breadth and the direct multilevel coupling are major MAPPINGS
strengths.

%==============================================================================
\section{Dust, PAHs, and PDR physics}

\subsection{Where the MAPPINGS microphysics is stronger}

The MAPPINGS grain microphysics is detailed.  The local source tree carries
graphite and silicate size distributions, depletion bookkeeping, grain
charge balance, electron and ion sticking, grain photoelectric heating,
collisional grain cooling, PAH charge states, and optional thermal or
ionization-parameter-dependent grain destruction.  For the local exchange
of energy and charge between gas and a resolved grain population, this is
more microscopic than the current \MoCHII{} Bakes--Tielens heating
prescription.

In \Mappings{}, this layer is connected to the P6/P7 zone integration and
to \code{dusttemp}, which computes a Draine--Li stochastic temperature
distribution and an IR photon field.  It is therefore a complete 1D dust
emission path, not merely a library capability.  The distinction below is
specific to the public \Mthree{} three-dimensional branch.

\subsection{Where the three-dimensional \texorpdfstring{M$^3$}{M3} dust path is limited}

The public \code{montph7\_dust} branch includes dust in the extinction and
terminates a packet when the selected interaction is grain absorption.
The dust data contain absorption, scattering, and asymmetry information,
but the examined 3D branch computes an absorption fraction and does not
sample a new dust-scattering direction.  The dedicated \code{dusttemp}
routine is called from the one-dimensional \code{photo6/7} family, not
from the public 3D driver.  Consequently, the rich MAPPINGS grain
microphysics should not be conflated with a complete 3D dust radiative
transfer and emergent-IR calculation.

This does not imply that dust is absent from \Mthree{}: dust opacity,
depletion, grain heating/cooling, and PAH processes affect the gas.  It
means that the public path does not presently expose a validated chain
equivalent to \MoCHII{}'s dust scattering, absorbed-energy ledger,
stochastic emission, and observer maps.

\subsection{The MoCHII dust chain}

\MoCHII{} couples dust to the same FUV/EUV radiation field that sets the
gas state:

\begin{enumerate}[leftmargin=2em]
\item Gas and dust compete for ionizing photons using bin-dependent
  absorption cross sections.  The dusty Str\"omgren gate reproduces the
  one-dimensional radius to $+0.69\%$ and the classic radius reduction
  $R/R_0=0.735$.
\item With scattering enabled, the extinction includes the dust scattering
  coefficient, a forced first event is sampled, and the packet follows an
  HG direction with an albedo weight.  The dusty-sphere result lies between
  the pure-absorption and extinction-as-absorption bounds.
\item The live dust model updates the dust abundance from the computed
  neutral fraction.  A separate PAH survival fraction permits destruction
  inside ionized gas and produces a PAH-bright I-front ring.
\item The absorbed FUV/EUV power gives a leaf equilibrium mixture
  temperature and an energy-conserving IR spectrum.
\item SEDust supplies astrodust, DL07, or Zubko stochastic emission.  The
  final normalization enforces
  $\int L_\lambda\,d\lambda=\sum_{\rm leaf}H_{\rm dust}V$; the smoke test
  closes at unity and resolves the usual PAH bands.
\item Leaf SEDs at selected wavelengths generate spatially resolved 8, 24,
  70, and 160-$\mu$m maps.
\end{enumerate}

For dusty three-dimensional H\,\textsc{ii} regions, this end-to-end chain
is the clearest overall physical advantage of \MoCHII{}.

\subsection{PDR scope}

\MoCHII{} extends the band below 13.6 eV so that low-threshold metals and
dust see the FUV field beyond the H I-front.  Optional switches include
electrons from the metal cascade, metal photoheating, Bakes--Tielens grain
photoelectric heating plus recombination cooling, and neutral-H impact
cooling in [C\,\textsc{ii}] 158 $\mu$m and [O\,\textsc{i}] 63 $\mu$m.
This is an atomic PDR skin, not a complete PDR code.  H$_2$/CO chemistry,
molecular self-shielding, molecular cooling, line transfer, and gas--dust
collisional coupling remain outside the model.

The MAPPINGS trees likewise contain more grain and general plasma physics
than a minimal H\,\textsc{ii} solver, but the present comparison did not
identify a full molecular H$_2$/CO network in either the 5.2.1
photoionization path or the public \Mthree{} 3D driver.  None of the three
codes should be selected as a substitute for a dedicated molecular PDR
code without additional chemistry.

%==============================================================================
\section{Emission lines, continua, and synthetic observations}

\subsection{Line inventories and emissivity accuracy}

The MAPPINGS lineage was designed to predict a very broad line spectrum.
The \Mthree{} cooling and output path includes resonance,
intercombination, forbidden, fine-structure, hydrogenic, helium, heavy
recombination, and extensive iron-line data.  Its five-level and
multilevel populations are already part of the thermal solution.  This
makes \Mthree{} attractive for integrated diagnostic grids and hard
nebulae for which a large fraction of the periodic table and many ion
stages matter.

\Mappings{} is the strongest version of that argument when 1D geometry is
acceptable: the full preference set covers 30 elements, and the
deterministic zone solver can afford a detailed spectrum without Monte
Carlo sampling noise.  \Mthree{} trades some parent-code recency and 1D
special modes for arbitrary Cartesian geometry.

\MoCHII{} has a narrower but highly auditable diagnostic layer.  Every
registered ion with a Tier-2 file is solved by the same partial-pivot
statistical-equilibrium code using fitted CHIANTI v11 effective collision
strengths and explicit $A$-values.  H I and He II recombination lines use
Storey--Hummer case-A/B tables.  The nebular continuum includes free-bound
plus free-free and two-photon components.  On an identical state, the
Fortran line luminosities match the independent Python solver to at least
four significant figures.  The limitation is coverage, not the internal
emissivity algebra.

\subsection{Images and observing geometry}

\MoCHII{} writes leaf state and line-emissivity blocks to HDF5/FITS.  Its
Python reader deposits each AMR leaf with exact pixel-area overlap, so
surface-brightness integration recovers the three-dimensional luminosity
to machine precision.  It also supports slab cuts and EM-, electron-,
hydrogen-, or volume-weighted field maps.

For the transported FUV/EUV field, an extra converged-state pass applies
observer peel-off.  Stellar and diffuse emission contribute direct
$e^{-\tau}$ terms, every dust interaction contributes an HG-weighted
scattered term, and the optical-depth walk uses the same gas and dust
opacity as the rate calculation.  Multiple observers, unattenuated direct
images, and bin-resolved cubes are available.  The direct thin-medium gate
matches the analytic image-integrated flux by $+0.15\%$ under the current
defaults (exact hydrogenic cross sections, threshold-aligned bins).

\Mthree{} outputs cell-resolved physical quantities and line emissivities,
and its published applications construct spatial line maps from them.
This is adequate for optically thin line projection.  The public driver
does not show a comparable general next-event estimator for arbitrary
observers, nor the separation of direct and dust-scattered continua.
Thus \Mthree{} is richer in the number of intrinsic lines, while
\MoCHII{} is richer in the chain that turns a three-dimensional state
into calibrated synthetic observations.

%==============================================================================
\section{Convergence, verification, and maturity}

\subsection{Convergence criteria}

\Mthree{} updates the radiation field, solves each illuminated cell, and
counts cells whose temperature and H$^+$ changes meet the selected
tolerances.  The run stops when a requested fraction of ionized cells has
converged or after the iteration limit.  Its public defaults use percent-
level cell changes and increase the packet count when convergence stalls.
This is intuitive, but the result is sensitive to how the ionized-cell
mask and required fraction are chosen.

\MoCHII{} reports both cell-max and volume-integrated pairs:
\begin{equation}
 \delta_x^{(V)} =
 \frac{\sum V\,|\Delta x_{\rm HII}|}{\sum V x_{\rm HII}},
 \qquad
 \delta_T^{(V)} =
 \frac{\sum V\nelec\,|\Delta\Te|}{\sum V\nelec\Te}.
 \label{eq:conv}
\end{equation}
The volume criterion suppresses front-cell jitter by the small volume of
the front and excludes temperature oscillations in nearly electron-free
cells.  In the FUV dusty test, cell maxima remained at 0.1 and 6.4 while
the volume measures reached stable Monte Carlo floors of
$2.7\times10^{-3}$ and $6.6\times10^{-3}$.  Thirty further iterations
changed the ionized volume by only 0.014\%.  This is a stronger convergence
diagnostic for turbulent, clumpy grids, provided the tolerance is set just
above the packet-noise floor.

\subsection{Validation evidence}

The MAPPINGS lineage has the stronger external maturity record.  \Mthree{}
inherits decades of MAPPINGS development and has been published against the HII20, HII40,
and PN150 Lexington/Meudon standards.  Its three-dimensional blister,
bipolar, and fractal applications demonstrate the importance of geometry
for spatially resolved line diagnostics \cite{jin2022,jin2022fractal}.

\MoCHII{} is newer, but its current repository contains unusually granular
gates:

\begin{itemize}[leftmargin=2em]
\item analytic H/He attenuation and rate integrals;
\item a Str\"omgren sphere against an exact one-dimensional reference;
\item AMR versus uniform and solution-driven refinement versus native high
  resolution;
\item H/He thermal balance against an independent Gauss--Seidel reference;
\item direct CHIANTI fit evaluation and PyNeb line-ratio checks;
\item MOCASSIN HII20/HII40 state and line comparisons;
\item case-A/case-B diffuse-field bracketing;
\item dust absorption/scattering brackets and IR energy closure;
\item direct peel-off flux against an analytic observer solution; and
\item exact three-dimensional emissivity-to-map luminosity conservation.
\end{itemize}

These gates provide strong internal evidence and make regressions easy to
localize.  They do not yet replace the independent user base and long
publication history of MAPPINGS.  The correct conclusion is that
the MAPPINGS lineage is more mature externally, while \MoCHII{} is more explicit and
fine-grained in its current numerical audit trail.

\Mappings{} does not have a Monte Carlo convergence problem: zone step
control and the global shock/precursor iteration dominate instead.  Its
decades of use and broad model suite are a major maturity advantage, but
the inspected 5.2.1 tree does not provide a gate hierarchy as explicit as
the current \MoCHII{} test directories.  A controlled local benchmark is
still required before attributing any line difference to geometry rather
than to the 5.1.13-versus-5.2.1 atomic and solver divergence.

%==============================================================================
\section{Software architecture and maintainability}

\MoCHII{} uses Fortran modules, allocatable leaf arrays, a single namelist,
MPI-3 shared memory, HDF5/FITS sections, and a data-driven species
registry.  New atomic products are generated offline by Python tools and
carry machine-readable provenance.  A new ion normally requires no new
solver code.  The transport, ion balance, thermal balance, lines, dust,
and output layers are separated modules.  This architecture substantially
reduces the cost and risk of adding physics.

The public \Mthree{} tree retains the Fortran 77 MAPPINGS architecture:
large common blocks, compile-time maxima, interactive input, and many
global arrays.  Several historical \code{montph7\_*} variants duplicate
large parts of the transport driver.  This structure has the advantage of
preserving a mature code base and its full menu of MAPPINGS models, but it
makes a change to the 3D algorithm harder to propagate and test.  The
available Sphinx input/output documentation is sparse, so reliable use
often requires reading the source and reproducing the authors' setup.

The architectural verdict is therefore clear: \MoCHII{} is easier to
extend, audit, automate, and connect to simulation data.  \Mthree{}'s
value lies primarily in the depth of the inherited physical library, not
in modern software structure.

\Mappings{} shares the same common-block and interactive heritage, but its
organization is cleaner than the collection of duplicated \Mthree{}
\code{montph7\_*} experiments and it has a single public \code{map52}
model selector.  Runtime preference files make the 16- or 30-element atom
sets configurable without recompilation.  Nevertheless, automated model
grids require scripting interactive input and parsing many ASCII outputs;
the current documentation overview explicitly remains incomplete.  For
maintainability the ordering is \MoCHII{} first, \Mappings{} second, and
the experimental \Mthree{} 3D fork third.

%==============================================================================
\section{Suitability by science case}

\begin{longtable}{P{0.42\textwidth}P{0.20\textwidth}P{0.30\textwidth}}
\caption{Recommended code by present science requirement.}
\label{tab:science}\\
\toprule
Science requirement & Preferred code & Reason \\
\midrule
\endfirsthead
\multicolumn{3}{l}{\small\itshape Table~\ref{tab:science} continued.}\\
\toprule
Science requirement & Preferred code & Reason \\
\midrule
\endhead
\bottomrule
\endfoot
Clumpy or porous H\,\textsc{ii} region from a TNG/RAMSES snapshot &
\MoCHII{} & AMR, front refinement, and simulation-oriented I/O. \\

Ionization-front position, shadows, and clump skins & \MoCHII{} &
Adaptive resolution and the analytic absorption estimator. \\

Dusty Str\"omgren structure and UV attenuation & \MoCHII{} & Gas--dust
competition and validated HG scattering in the 3D band. \\

PAH rings, IR SEDs, or 8--160-$\mu$m maps & \MoCHII{} & SEDust,
leaf dust emissivities, and exact absorbed-energy normalization. \\

Spatially resolved IFU line maps & Usually \MoCHII{} & AMR emissivities
and flux-conserving projection; verify that the required ions exist. \\

Very broad integrated optical/UV/IR line inventory & \Mappings{} & Current
full 30-element MAPPINGS preference set and deterministic detailed spectrum. \\

Compact high-density nebula & \Mappings{} if 1D; \Mthree{} if 3D & Both
feed density-dependent multilevel cooling back into $\Te$; geometry decides. \\

Hard PN-like ionizing spectrum & \Mappings{} if 1D; \Mthree{} if 3D & More
high charge states; \Mthree{} additionally has published PN150 3D validation. \\

Non-Maxwellian $\kappa$ electron experiment & \Mappings{} or \Mthree{} &
Both MAPPINGS-derived trees have native $\kappa$ rate support. \\

One-dimensional radiation-pressure-confined nebula & \Mappings{} &
Isobaric photoionization can include the accumulated radiation pressure. \\

Finite source lifetime or source-off ionization history in 1D & \Mappings{}
& P6/P7 connect finite-age modes to the time-dependent ion solver. \\

Steady 1D radiative MHD shock and precursor & \Mappings{} & S5 couples
adaptive post-shock cooling to an iterated auto-preionization solution. \\

CIE, NEQ, or PIE cooling-curve grid & \Mappings{} & Dedicated CC, NC, and
PP drivers and a broad general-plasma ion inventory. \\

Controlled test of modern CHIANTI v11 coefficients & \MoCHII{} &
Versioned fitted data and direct fit-error gates. \\

Three-dimensional time-dependent ionization & None & Both public 3D
drivers are equilibrium solvers; \Mappings{} time dependence is 1D. \\

Three-dimensional radiative shock & None & \Mappings{} has a strong 1D
shock mode, but it is not coupled to either 3D transport architecture. \\

Molecular H$_2$/CO PDR & None & All three need a dedicated molecular
chemistry and line-transfer layer. \\

Ly$\alpha$ or metal resonance-line profiles & None & Requires
frequency-redistributing resonant transfer, e.g.\ LaRT. \\
\end{longtable}

%==============================================================================
\section{Recommended development path for MoCHII}

The comparison suggests a focused roadmap rather than wholesale
replacement of either code.

\begin{enumerate}[leftmargin=2em]
\item \textbf{Put density dependence into the thermal cooling.}  Reuse the
  generic $n$-level solver in the thermal loop for dominant coolants, or fit
  a suppression factor versus $(\Te,\nelec)$ around each critical density.
  This closes the most consequential gap relative to MAPPINGS.
\item \textbf{Extend the registry upward in ion stage.}  Add the ions needed
  for PN and hard-field applications before adding marginal new low stages.
  The existing data-operation design is well suited to this.
\item \textbf{Add Na, Al, and Ni only when driven by a diagnostic need.}
  Their MAPPINGS implementations provide useful reference behavior, while
  the \MoCHII{} pipeline should retain explicit provenance and ion-by-ion
  gates.
\item \textbf{Retain the present transport architecture.}  The octree,
  analytic estimator, Cartesian DDA, threshold-aligned bins, and peel-off
  layer are major advantages and should not be replaced by the legacy
  \Mthree{} packet synchronization model.
\item \textbf{Use \Mappings{} as the differential microphysics oracle.}
  Construct one-zone tests at fixed $(\Te,\nelec,J_\nu)$ and compare ion
  fractions, heating/cooling terms, and individual line emissivities before
  comparing full 3D nebulae.  Use \Mthree{} only as the second-stage test
  of how that inherited microphysics behaves in arbitrary geometry.
\item \textbf{Perform a controlled common-physics benchmark.}  Use the same
  uniform grid, source SED, abundances, dust-off setting, frequency range,
  and atomic subset.  Compare $J_\nu$, photo-rates, ion fractions, $\Te$,
  cooling budgets, and line emissivities separately.  Only then compare
  total runtime and MPI scaling.
\end{enumerate}

The most valuable hybrid would combine \Mappings{}-like density-dependent
microphysics with the current \MoCHII{} geometry, transport, dust, and
observer layers.  This would preserve the strongest feature of each code
without inheriting the principal numerical and software limitations of the
other.

%==============================================================================
\section{Final assessment}

\begin{center}
\begin{tabular}{P{0.46\textwidth}P{0.20\textwidth}P{0.20\textwidth}}
\toprule
Capability & Stronger now & Confidence \\
\midrule
Three-dimensional transport estimator & \MoCHII{} & High \\
Adaptive spatial resolution & \MoCHII{} & High \\
Dust scattering and IR/PAH observables & \MoCHII{} & High \\
Observer imaging and data products & \MoCHII{} & High \\
Atomic-data provenance and fit auditability & \MoCHII{} & High \\
Element and ion-stage breadth & \Mappings{} & High \\
Density-dependent thermal line cooling & \Mappings{}/\Mthree{} & High \\
Hard PN/general plasma coverage & \Mappings{} & High \\
Finite-age 1D photoionization & \Mappings{} & High \\
Radiative MHD shock and precursor & \Mappings{} & High \\
Noise-free 1D parameter-grid throughput & \Mappings{} & High \\
External validation history & MAPPINGS lineage & High \\
Software extensibility and automation & \MoCHII{} & High \\
Absolute line accuracy for every ion & Case dependent & Requires a
common-data benchmark \\
Massively distributed memory ceiling & Architecture dependent & Requires
a scaling experiment \\
\bottomrule
\end{tabular}
\end{center}

For the principal intended application--dusty, structured, spatially
resolved H\,\textsc{ii} regions--\MoCHII{} is the better current platform.
For broad integrated emission-line diagnostics, dense nebulae, harder
ionization conditions, finite-age 1D photoionization, and radiative shocks,
\Mappings{} is the more complete physical reference.  \Mthree{} becomes
preferable only when those MAPPINGS-like line and cooling strengths must be
combined with arbitrary three-dimensional Cartesian geometry.  None of
these judgments is universal: first decide whether geometry, time history,
shock dynamics, dust observables, density-dependent cooling, or ion-stage
breadth is the controlling requirement.

%==============================================================================
\section{Hard-spectrum (PN and AGN) extension plan}
\label{sec:hardplan}

This section records the staged plan that follows from the comparison
(kept as \code{docs/HARD\_SPECTRUM\_PLAN.md}, drafted 2026-07-18, not yet
started).  It uses the same staged-gate discipline as \code{docs/PLAN.md}.

\subsection{Scope and the two tracks}

``Hard spectrum'' splits into two physically distinct regimes, kept as
separate tracks because they need different physics:

\begin{itemize}[leftmargin=2em]
\item \textbf{PN track} --- thermal central stars,
  $T_{\rm eff}\sim75$--200 kK.  A 150 kK blackbody has $kT=12.9$ eV; its
  photon flux at 300 eV sits $\sim$23\,$kT$ out on the Wien tail
  ($e^{-23}\sim10^{-10}$ before the phase-space factors), and the lowest
  K-shell edge of an abundant metal (C, 290 eV) is never reached with
  significant flux.  The MOCASSIN PN150 benchmark itself integrates its
  band only to \code{nuMax} $=24$ Ryd $=326$ eV.  The PN track therefore
  needs \emph{no inner-shell/Auger physics and no Compton terms} --- it
  is band headroom, higher ion stages, the \HeII{} diffuse channels, and
  (at PN densities) density-dependent cooling.
\item \textbf{AGN track} --- a power-law/big-bump continuum extending to
  keV X-rays.  Above $\sim$0.3 keV inner-shell absorption dominates the
  metal opacity and Auger multi-electron ejection couples non-adjacent
  ion stages; near $\sim$0.5--1 keV the Thomson cross section
  ($6.65\times10^{-25}$ cm$^2$) overtakes the steeply falling
  ($\sigma\propto E^{-3}$) photoabsorption of the residual neutrals, and
  Compton heating enters the energy budget.  The first increment caps
  the band at $\sim$500 eV, where photoabsorption still dominates, and
  defers the Compton terms (Section~\ref{sec:harddeferred}).
\end{itemize}

Current baseline (v1.00-dev, 2026-07-18): band 13.6--100 eV (+FUV
option); registry stages C 4 / N 3 / O 3 / Ne 3 / S 4 / Ar 4 / Mg 3 /
Fe 4 / Si 5 / Cl 5 / Ca 5, with the compile-time stage ceiling
\code{MAX\_ST} $=6$ in \code{species\_mod} --- the PN targets below fit
it exactly, so no framework change is needed there; \code{te\_max}
$=5\times10^4$ K and the Tier-1 cooling fits are generated over
$10^3$--$10^5$ K; $\bar g_{\rm ff}$ is tabulated for $Z_{\rm ion}=1$--4
(\code{cooling\_mod}, parameter \code{NGZ}); secondary ionization
(\code{use\_sec\_ion}, Shull--van Steenberg 1985) is implemented but
default off; the Owen-scrambled Sobol launch is available.  The \HeII{}
Storey--Hummer lines (case A and B), the \HeI{} Porter lines, and the
diffuse \HeIII{} ground channel already cover the \HeII{} 4686
diagnostic side, and the aligned-edge band builder
(\code{ion\_align\_edges}, default on) with its graceful
threshold-dropping logic is exactly the piece that makes many new
in-band edges safe.

\subsection{Reference assets (verified present)}

\begin{longtable}{P{0.44\textwidth}P{0.46\textwidth}}
\caption{Local reference assets for the hard-spectrum work.}
\label{tab:hardassets}\\
\toprule
Asset & Use \\
\midrule
\endfirsthead
\multicolumn{2}{l}{\small\itshape Table~\ref{tab:hardassets} continued.}\\
\toprule
Asset & Use \\
\midrule
\endhead
\bottomrule
\endfoot
MOCASSIN PN150 + PN75 benchmark inputs
(\code{benchmarks/gas/\{PN150,PN75\}}) & primary 3D Monte Carlo gates;
run MOCASSIN itself \emph{converged} (the HII40 lesson).  PN150 input
read directly: 150 kK blackbody, $n_{\rm H}=3000$ cm$^{-3}$,
$R_{\rm in}=10^{17}$, $R_{\rm out}=4.07\times10^{17}$ cm, band to
24 Ryd, abundances He 0.10, C $3\times10^{-4}$, N $10^{-4}$,
O $6\times10^{-4}$, Ne $1.5\times10^{-4}$, Mg $3\times10^{-5}$,
Si $3\times10^{-5}$, S $1.5\times10^{-5}$ (\code{abunPN150.in}) ---
note Mg and Si ARE in the composition, which drives their stage targets
below. \\
Cloudy c23.01 test suite (\code{tsuite/auto/pn\_paris*.in},
\code{nlr\_paris.in}, \code{agn\_lex00\_u\{0,1,m1\}.in}) & 1D cross-code
gates: Meudon PN, AGN NLR, and the Lexington-2000 X-ray slabs.  The
\code{agn\_lex00} inputs read directly: $n_{\rm H}=10^5$ cm$^{-3}$,
$\phi({\rm H})=10^{15.477}$ cm$^{-2}$s$^{-1}$, stop column
$10^{16}$ cm$^{-2}$, and a \emph{tabulated} SED (Cloudy
\code{interpolate/continue} points) --- so \MoCHII{} reproduces the
illumination exactly through the existing \code{spectrum\_type} file
route, without waiting for the AGN preset. \\
Cloudy AGN SED (\code{source/parse\_agn.cpp}) & the standard AGN
continuum parameterization (big-bump $T$, $\alpha_{\rm ox}$ anchored at
2500 \AA--2 keV, $\alpha_{\text{uv}}$, $\alpha_{\rm x}$) for a \MoCHII{}
preset. \\
Kaastra \& Mewe (1993) Auger yields
(\code{data/mewe\_nelectron.dat}, \code{mewe\_fluor.dat}) & electrons
ejected per inner-shell vacancy, shell-resolved, machine-readable
provenance for the Auger stage. \\
Verner \& Yakovlev (1995) shell-resolved cross sections
(\code{sigma\_vy95} already in \code{photo\_xsec.f90}) & inner-shell
photoionization above the K/L edges. \\
CHIANTI v11.0.2 + Badnell 2023 RR/DR (in-tree) & levels/collision
strengths and recombination for every new ion stage (coverage
verified). \\
\end{longtable}

\subsection{HS1: band headroom}

Mostly validation; the code changes are one-line parameters.

\begin{enumerate}[leftmargin=2em]
\item \code{eion\_max} (define.f90, default 100 eV) raised to
  300--500 eV for hard runs.  This matters even for stages the registry
  ALREADY tracks: the Ne\,\textsc{iv}$\to$\textsc{v} threshold is
  97.1 eV --- just under the present cap --- so today's band cannot
  photoionize to Ne\,\textsc{v} at all, and the new O/Ne thresholds
  (113.9, 126.2 eV) lie beyond it.  Cross-section validity is not the
  issue (the Verner et al.\ fits extend to keV energies and \HI/\HeII{}
  are exact hydrogenic); the work is verifying the pieces that ASSUME
  the band shape: the aligned-edge builder with $\sim$30--40 in-band
  thresholds once HS2 lands (recommend \code{nnu\_ion} $\geq 64$, 128
  clean; the largest-remainder allocator and the metal-edge dropping
  logic already degrade gracefully), the FUV+ionizing two-segment
  bookkeeping, the diffuse \code{ion\_bin\_of} edge search, and the
  secondary-ionization band integrals whose $E_0>40$ eV branch becomes
  the COMMON case on a hard field.
\item $\bar g_{\rm ff}$: raise \code{NGZ} from 4 to 8 in
  \code{cooling\_mod} (the van Hoof table already spans the required
  $\gamma^2$ range; the metal free-free sum then covers the new
  stages).
\item \code{te\_max}: allow $10^5$--$10^6$ K for X-ray-heated zones.
  Any ion used above $10^5$ K needs its Tier-1 fit regenerated over the
  wider range (a flag in \code{chianti\_cooling.py}, not new code); the
  bisection itself is range-agnostic.
\item Recommended hard-field defaults, documented in the example
  inputs: \code{use\_sec\_ion=.true.}\ (a 150 kK field makes
  $E_0>40$ eV photoelectrons routine, and X-ray fields more so) and
  \code{launch\_sequence='sobol'} (the shielded, partially neutral
  zones where secondaries act are exactly where launch noise hurts
  most).
\end{enumerate}

\emph{Gate HS1:} the G0 pattern on the widened band --- fixed-state
analytic attenuation with a 150 kK blackbody, $\Gamma_{\rm HI/HeI/HeII}$
vs the analytic integrals bin by bin; $Q(>13.6)/Q(>24.6)/Q(>54.4)$
photon-rate ratios against exact Planck integrals to $<0.1\%$ with
aligned edges; and a spot check that the tabulated VFKY96 metal cross
sections at 300--500 eV match the published fit values (guards against
table-range assumptions in \code{photo\_xsec}).

\subsection{HS2: registry upward in ion stage}

A data operation through the existing fitting pipeline:

\begin{longtable}{P{0.10\textwidth}P{0.12\textwidth}P{0.28\textwidth}P{0.36\textwidth}}
\caption{Target stages and the diagnostics they unlock.}
\label{tab:hardions}\\
\toprule
El. & now$\to$target & new thresholds [eV] & key new lines \\
\midrule
\endfirsthead
\multicolumn{4}{l}{\small\itshape Table~\ref{tab:hardions} continued.}\\
\toprule
El. & now$\to$target & new thresholds [eV] & key new lines \\
\midrule
\endhead
\bottomrule
\endfoot
C & 4$\to$5 & C\,\textsc{iv}$\to$\textsc{v} 64.5 & C\,\textsc{iv}
1548/1550 (resonance --- optically thin caveat) \\
N & 3$\to$5 & 47.4, 77.5 & N\,\textsc{iv}] 1486, N\,\textsc{v} 1240 \\
O & 3$\to$6 & 54.9, 77.4, 113.9 & [O\,\textsc{iv}] 25.9 $\mu$m,
O\,\textsc{iv}] 1400, O\,\textsc{v}] 1218 \\
Ne & 3$\to$6 & 63.4, 97.1, 126.2 & [Ne\,\textsc{iv}] 2422/2425,
[Ne\,\textsc{v}] 3346/3426 + 14.3/24.3 $\mu$m, [Ne\,\textsc{vi}]
7.65 $\mu$m \\
S & 4$\to$6 & 47.2, 72.6 & [S\,\textsc{iv}] 10.5 $\mu$m (new $n$-level
file) \\
Ar & 4$\to$6 & 59.8, 75.0 & [Ar\,\textsc{v}] 7.90/13.1 $\mu$m,
[Ar\,\textsc{vi}] 4.53 $\mu$m \\
Mg & 3$\to$6 & 80.1, 109.3, 141.3 & stage tracking for the PN150
composition (Mg $3\times10^{-5}$); [Mg\,\textsc{iv}] 4.49 $\mu$m,
[Mg\,\textsc{v}] 5.61 $\mu$m as data permit \\
Si & 5$\to$6 & Si\,\textsc{v}$\to$\textsc{vi} 166.8 & stage tracking
for the PN150 composition (Si $3\times10^{-5}$) \\
\end{longtable}

Fe/Cl/Ca hold (extend only on diagnostic demand; coronal
[Fe\,\textsc{vii}+] deferred).  All targets fit the existing
\code{MAX\_ST} $=6$ ceiling exactly; \code{MAXLEV} $=40$ /
\code{MAXTR} $=600$ and the 600-line \code{lines\_mod} arrays already
accommodate the new $n$-level models.  Work items, all through
\code{tools/fitting/}:

\begin{enumerate}[leftmargin=2em]
\item \textbf{Element files} (\code{make\_element\_data.py}): VFKY96
  \code{PHOTO} rows exist in phfit2 for every target stage; Voronov CI
  covers all stages; Badnell 2023 RR+DR cover these isoelectronic
  sequences in the in-tree \code{clist\_K} files (with the CHIANTI
  \code{RR2}/\code{DR2} fallback forms already supported where a fit is
  absent).  \textbf{Charge exchange} above stage $\sim$4 is not in the
  Kingdon--Ferland fits: adopt the fast Dalgarno-type
  $\sim1.92\times10^{-9}z$ cm$^3$s$^{-1}$ recombination-only scaling
  used by Cloudy, with an explicit provenance note in each element
  file, or document neglect --- decide once, apply uniformly (CX
  matters little in the hot interior where these stages live, but the
  decision must be recorded).
\item \textbf{Tier-1 cooling fits} (\code{chianti\_cooling.py}) for
  each new stage, fit range matched to the chosen \code{te\_max}.
\item \textbf{Tier-2 $n$-level files} for the diagnostic ions
  ([Ne\,\textsc{iv}], [Ne\,\textsc{v}], [O\,\textsc{iv}],
  [S\,\textsc{iv}], [Ar\,\textsc{v}], [Ar\,\textsc{vi}], N\,\textsc{iv},
  N\,\textsc{v}, C\,\textsc{iv}); CHIANTI v11 carries level and
  collision data for all of them.  Ions without CHIANTI level data are
  tracked without lines --- the established Ar\,\textsc{ii} pattern.
\item \textbf{PN abundance example} (\code{examples/}): the
  \code{abunPN150.in} values verbatim, so the PN150 gate input is
  reproducible from the repository alone.
\end{enumerate}

\emph{Gate HS2:} PyNeb ratio checks ion by ion at the established
database-spread tolerances (sub-percent for well-determined optical
pairs, up to $\sim$5--10\% for the known compilation-spread classes):
the [Ne\,\textsc{v}] 3426/24.3 $\mu$m temperature pair, the
[Ne\,\textsc{iv}] 2422/2425 and [Ar\,\textsc{iv}] 4711/4740 density
pairs, [O\,\textsc{iv}] 25.9 $\mu$m against its A-value limit; then an
end-to-end hard-field smoke showing every new line in
\code{\_lines.txt} with physically sane ratios.

\subsection{HS3: He II excited-channel diffuse photons}

The PN \HeIII{} zone is large, and case-B \HeII{} recombinations emit
ionizing photons beyond the ground continuum already carried
(\code{diffuse\_mod} channel 3).  The hydrogenic $n=2$ cascade of
\HeII{} produces three components, every one of which ionizes \HI{}:

\begin{itemize}[leftmargin=2em]
\item \HeII{} Ly$\alpha$ at 40.8 eV --- ionizes \HI{} AND \HeI{}, the
  dominant piece;
\item the \HeII{} Balmer continuum --- its threshold is
  $Z^2\,{\rm Ry}/n^2 = 13.60$ eV, i.e.\ exactly at the \HI{} edge; the
  aligned-grid \code{ion\_bin\_of} binary search (the 2026-07-18 bug
  fix) is precisely what makes this land in the correct bin;
\item the \HeII{} two-photon continuum (0--40.8 eV), of which the
  $>13.6$ eV part ionizes --- sampled like the \HeI{} two-photon branch
  with its ionizing probability precomputed.
\end{itemize}

Implementation is the \code{hei\_diffuse} pattern verbatim: a fifth
channel in \code{dif\_ch} weighted by the case-B \HeIII{}
recombination rate in each leaf, branch fractions from the hydrogenic
$2{\rm s}/2{\rm p}$ cascade shares (Osterbrock \& Ferland tables),
energies fixed (Ly$\alpha$), threshold-plus-exponential (Balmer
continuum), or flat-sampled (two-photon).  The quasi-random diffuse
launch needs NO new dimensions --- the channel pick d7 simply spans
five values.  Regime note: at HII40 the analogous \HeI{} channel moved
$\Te$ by only $\sim$1\%, but there the \HeIII{} zone was a small core;
in a 150 kK PN the \HeIII{} region fills much of the nebula and this
becomes a first-order energy-budget term.

\emph{Gate HS3:} channel energy conservation (emitted diffuse
luminosity equals the case-B $n\geq2$ recombination power of the
\HeIII{} zone); the PN150 He structure and $\Te$ recorded with the
channel off/on, attributed as one lever.

\subsection{HS4: PN validation gates}

\textbf{MOCASSIN PN150/PN75}, with the full-quality protocol that
resolved HII40: MOCASSIN re-run CONVERGED (its stored reduced settings
--- $10^6$ photons, 10 iterations, \code{convLimit} 0.05, $13^3$ ---
are a moving target exactly as the HII20/40 baselines were), and
\MoCHII{} run with \code{conv\_crit='vol'}, \code{nnu\_ion} 64--128
aligned, $\geq4\times10^7$ packets, \code{launch\_sequence='sobol'}.
The PN150 setup is fully specified by the benchmark input quoted in
Table~\ref{tab:hardassets} (150 kK, $n_{\rm H}=3000$ cm$^{-3}$, shell
$10^{17}$--$4.07\times10^{17}$ cm).  Compare, in order of diagnostic
power:

\begin{enumerate}[leftmargin=2em]
\item $\Te$(\HII{}) and the radial $\Te$ profile (PN Te $\sim$12--18 kK
  --- also exercises the hotter thermal range);
\item the He\,\textsc{i}/\textsc{ii}/\textsc{iii} structure --- the
  \HeII{}-edge (54.4 eV) photon budget is the PN analog of the He edge
  discretization that dominated HII40, so the aligned grid must show
  its value here;
\item the new high-stage fractions O\,\textsc{iii/iv/v},
  Ne\,\textsc{iii/iv/v}, N\,\textsc{iii/iv};
\item the line table: $L$(H$\beta$), \HeII{} 4686/H$\beta$,
  [O\,\textsc{iii}] 5007/4363, [Ne\,\textsc{v}] 3426,
  [Ne\,\textsc{iii}] 15.5 $\mu$m, [O\,\textsc{iv}] 25.9 $\mu$m,
  [N\,\textsc{ii}] 6584, [S\,\textsc{iv}] 10.5 $\mu$m,
  [Ar\,\textsc{v}].
\end{enumerate}

\textbf{Cloudy \code{pn\_paris}} as the 1D cross-check with identical
SED, abundances, density, and dust off: zone-by-zone $\Te$ and ion
fractions along the radius plus $\sim$20 integrated ratios.  Because
Cloudy solves density-dependent level populations in its thermal loop,
the PN150 residual BEFORE stage HS7 directly measures how much the
low-density Tier-1 fits miss at $n_{\rm H}=3000$ cm$^{-3}$ --- that
number is the input to HS7, which is why HS4 deliberately precedes it.

\emph{Acceptance:} land inside (or bracket) the published Lexington
PN150 cross-code spread, with each residual attributed one lever at a
time (binning, diffuse channels, density-dependent cooling, data
vintage) exactly as the HII40 campaign did.

\subsection{HS5: AGN SED input}

An \code{ion\_spectrum='agn'} preset implementing the Cloudy continuum
(\code{parse\_agn.cpp}; Mathews--Ferland form):
\begin{equation}
 f_\nu \;=\; \nu^{\alpha_{\text{uv}}}\,
   e^{-h\nu/kT_{\rm bump}}\, e^{-kT_{\rm IR}/h\nu}
   \;+\; a\,\nu^{\alpha_{\rm x}} ,
 \label{eq:agnsed}
\end{equation}
with the X-ray normalization $a$ fixed by $\alpha_{\rm ox}$ through the
2500 \AA--2 keV energy ratio (the 403.3 factor in the Cloudy source)
and the standard defaults $\alpha_{\text{uv}}=-0.5$,
$\alpha_{\rm x}=-1$, IR cutoff at $kT_{\rm IR}=0.01$ Ryd.  It enters
\code{ion\_band\_mod} exactly like the Planck function --- a spectral
density integrated on the 32-point sub-grid in each band bin --- and is
normalized by the existing derive-or-rescale luminosity rules.  The
tabulated-file route (\code{spectrum\_type='per\_ev'} etc.) remains the
reference path and already reproduces the \code{agn\_lex00} gate SEDs
verbatim (Table~\ref{tab:hardassets}); the preset adds convenience and
reproducibility for parameter studies.  Also add an
ionization-parameter helper to the rank-0 log,
$U = Q_{\rm ion}/(4\pi r^2 n_{\rm H} c)$ evaluated at a reference
radius, so slab setups map onto Cloudy inputs without hand conversion.

\emph{Gate HS5:} the bin-integrated preset spectrum against Cloudy
\code{save continuum} output for identical parameters (shape agreement
to $\sim$1\% per bin); $Q$-ratio integrals exact; the derive/rescale
bookkeeping re-verified on the preset (the ISRF-preset test pattern).

\subsection{HS6: inner-shell absorption and Auger (AGN track)}

The one genuine framework extension of the plan, in four parts:

\begin{enumerate}[leftmargin=2em]
\item \textbf{Shell-resolved opacity.}  Above each K/L edge the total
  cross section becomes the VFKY96 outer-shell fit plus the Verner \&
  Yakovlev (1995) inner-shell partials,
  $\sigma_{\rm tot}(E)=\sigma_{\rm outer}(E)+\sum_s\sigma_s(E)$.  The
  \code{sigma\_vy95} evaluator and the \code{PHOTO2} element-file rows
  already exist (the Cl path); the extension is carrying several shells
  with their edges in the aligned-band threshold list.
\item \textbf{Auger yields.}  Kaastra \& Mewe (1993) electrons per
  vacancy, shell-resolved, read from the Cloudy data file format
  (\code{mewe\_nelectron.dat}: element, ion, shell, yield
  distribution): an inner-shell ionization of stage $i$ lands at stage
  $i+n_{\rm e}({\rm shell})$ with the tabulated probability
  distribution over $n_{\rm e}$.  Fluorescence (K$\alpha$ photon
  emission instead of one Auger electron, \code{mewe\_fluor.dat}) is
  a small branch for low-$Z$ elements; the first increment deposits it
  as local heat with a documented note rather than transporting the
  line photon.
\item \textbf{Cascade generalization.}  \code{species\_fractions}
  currently solves the strict adjacent-stage chain as a product form.
  Auger couples stage $i$ to $i+2,\,i+3,\dots$, so each element's
  stationary balance becomes a small $n_{\rm stage}\times n_{\rm stage}$
  matrix problem ($M\mathbf{x}=0$ with $\sum_j x_j=1$); the
  partial-pivot elimination already used by \code{nlevel\_mod} is
  reused directly.  The chain path is KEPT as the default and the
  matrix path activates only when an inner-shell channel is present ---
  the off-path stays arithmetically identical, the house rule for every
  extension.
\item \textbf{Energetics.}  The Auger electron carries
  $E_0 = h\nu - E_{\rm th,shell} - \sum E_{\text{secondary vacancies}}$
  and is fast; it enters the EXISTING Shull--van Steenberg partition
  through the absorber-resolved excess-energy integrals with no new
  heating physics.  Compton remains excluded at $\leq$500 eV.
\end{enumerate}

\emph{Gate HS6:} the X-ray-illuminated slab against the Cloudy
\code{agn\_lex00\_u\{m1,0,1\}} set (three ionization parameters at
$n_{\rm H}=10^5$ cm$^{-3}$, tabulated SED, stop column
$10^{16}$ cm$^{-2}$): ion-fraction profiles of C/N/O/Ne through the
partially ionized zone, the $\Te$ profile, and the H/He recombination
lines, within the published Lexington-2000 cross-code spread.  A
secondary check: with Auger artificially disabled the slab must
reproduce the HS2-era solution (isolates the new coupling).

\subsection{HS7: density-dependent metal cooling in the thermal loop}

Required at PN densities.  The coolants that dominate at nebular
temperatures quench in a specific order because their critical
densities span five decades (approximate values at $10^4$ K):
[C\,\textsc{ii}] 158 $\mu$m $n_{\rm crit,e}\sim50$,
[N\,\textsc{ii}] 122 $\mu$m $\sim3\times10^2$,
[O\,\textsc{iii}] 88 $\mu$m $\sim5\times10^2$,
[O\,\textsc{iii}] 52 $\mu$m $\sim3.5\times10^3$,
[O\,\textsc{ii}] 3729 $\sim3\times10^3$,
[S\,\textsc{ii}] 6716 $\sim2\times10^3$,
[N\,\textsc{ii}] 6584 $\sim9\times10^4$,
[O\,\textsc{iii}] 5007 $\sim7\times10^5$.  At the PN150 density
($n_{\rm H}=3000$, $\nelec$ up to $\sim$7000 cm$^{-3}$ in the \HeIII{}
zone) the IR fine-structure lines are already partly quenched while the
optical forbidden lines are not --- the low-density Tier-1 fits
therefore OVERESTIMATE cooling and bias $\Te$ low, with the error
growing toward compact PN and NLR densities.

Primary design, keeping ``fitted offline, evaluated online'':

\begin{enumerate}[leftmargin=2em]
\item \code{chianti\_cooling.py} already solves the full $n$-level
  system to build its fits; extend it to emit
  $\Lambda(T,\nelec)$ tables on a log-$\nelec$ grid ($\sim$7 nodes,
  $10^1$--$10^7$ cm$^{-3}$) $\times$ the existing log-$T$ range, one
  block appended to each \code{cooling\_tier1\_*.txt} behind a
  format-version flag (old files keep working).
\item \code{species\_mod::metal\_cooling} evaluates
  bilinear-in-$(\log T,\log\nelec)$; the memory cost is trivial
  ($\sim$60$\times$7 values for each ion) and the evaluation stays a
  handful of multiply-adds inside the bisection --- the
  \code{species\_ne} cache pattern applies unchanged.
\item The lowest-density node MUST reproduce the current fits to fit
  accuracy (regression guard), so HII20/HII40 and every recorded gate
  move by $<0.1\%$.
\item Verification: an in-loop $n$-level solve for the dominant
  coolants on a small model, confirming table interpolation error
  $<1\%$ in $\Lambda$ over the used $(T,\nelec)$ range; the H-impact
  PDR cooling terms (\code{metal\_cooling\_H}) stay as they are (their
  two-level low-density form is correct in the PDR regime they serve).
\end{enumerate}

\emph{Gate HS7:} a two-zone analytic check (one low- and one
high-density zone against the exact $n$-level cooling); PN150 $\Te$
with the tables off/on (the residual measured at HS4 must shrink
accordingly); HII20/HII40 unchanged to $<0.1\%$.

\subsection{Deferred items}
\label{sec:harddeferred}

Recorded once so they are conscious exclusions, not oversights:
Compton heating/cooling and Thomson-scattering transport
($\gtrsim$0.5 keV; Cloudy is the reference implementation); Fe K
fluorescence photon transport (the yield enters HS6 as local heat);
coronal ions beyond $\sim$stage 8 ([Fe\,\textsc{vii}] 6087,
[Fe\,\textsc{x}] 6374 need Fe stages 8--11 and a hotter thermal range);
X-ray and UV resonance-line transfer (LaRT territory); grain
destruction/sputtering in hard radiation fields (the PAH survival
switches are a partial hook, not a destruction model); and time
dependence / finite source age, where MAPPINGS remains the 1D
reference.

\subsection{Suggested order}

\begin{center}\small
\begin{tabular}{P{0.06\textwidth}P{0.38\textwidth}P{0.20\textwidth}P{0.20\textwidth}}
\toprule
Order & Stage & Framework risk & Blocking \\
\midrule
1 & HS1 band headroom & low & gates HS2+ \\
2 & HS2 registry upward & none (data) & gates HS4 \\
3 & HS3 \HeII{} diffuse channels & low & sharpens HS4 \\
4 & HS4 PN gates (MOCASSIN + Cloudy) & none & PN track done \\
5 & HS7 density-dependent cooling & medium & PN accuracy \\
6 & HS5 AGN SED preset & low & gates HS6 \\
7 & HS6 inner-shell/Auger & high (cascade) & AGN track done \\
\bottomrule
\end{tabular}
\end{center}

HS7 can run in parallel with HS2--HS4 (independent code paths); the
first PN150 pass deliberately precedes HS7, so it measures how much of
the PN residual the density-dependent cooling must explain --- the same
one-lever-at-a-time attribution that worked for HII40.

%==============================================================================
\section{Feasibility: time dependence, cooling flows, and shocks in
\texorpdfstring{\MoCHII{}}{MoCHII}}
\label{sec:tdcs}

Section~\ref{sec:tdcs}\ addresses the three capabilities where the public
\Mthree{} 3D driver and \MoCHII{} both trail \Mappings{}: the connected
time-dependent ionization history of P6/P7, the 1D cooling flow, and the
S5 steady radiative shock with an iterated precursor
(\S\ref{sec:mapstrength}).  The question is not whether these are useful
--- they plainly are --- but how much each would cost to add to
\MoCHII{}, given its architecture.  The three answers differ sharply:
time dependence is a bounded project that the design accommodates
naturally, a cooling flow rides on top of it, and a radiative shock is a
separate one-dimensional radiation-hydrodynamic calculation that the
Monte Carlo transport engine does little to help.

\subsection{The architectural pivot}

\MoCHII{}'s iteration loop transports Monte Carlo packets, tallies the
cell radiation field, solves the \emph{equilibrium} ionization and thermal
balance in each cell (\code{solve\_ion\_cell}: a damped fixed point on
$\nelec$ with a bisection on $\Te$), refreshes the opacity, and repeats
until converged.  The state solver is therefore a \emph{root finder}:
given the rates it returns the steady state.  It integrates no
$\mathrm{d}n_i/\mathrm{d}t$, carries no velocity field, and solves no
fluid equations.  That single fact sets the cost of each addition: one
reuses the root finder almost unchanged, one adds a time integrator
beside it, and one needs what the code does not have at all.

\subsection{Time-dependent ionization: moderate, and the natural fit}

Almost everything a finite-lifetime calculation needs already exists.  The
transport, the photoionization, recombination, collisional-ionization, and
charge-exchange rate coefficients (already evaluated cell by cell), the
tally, the octree, and the diffuse field are all reusable without change.
The cell solver is already ``given the rates, find the state''; the new
piece is a sibling routine, ``given the rates, the old state, and
$\Delta t$, find the new state.''

What is genuinely new is threefold.  First, a stiff ordinary-differential
integrator for the ionic populations and, where the thermal time is not
short compared with the ionization time, the electron temperature --- a
backward-Euler or semi-implicit step in each cell, of the kind
\Mappings{}'s \code{timion} performs, rather than the present root find.
Second, an outer time loop with a step controller keyed to the atomic,
cooling, and photon-absorption times.  Third, a source light curve or a
finite on/off lifetime as input.

The design accommodates this cleanly through operator splitting.
\MoCHII{} already performs one transport pass for each iteration, so the
natural structure is an outer radiation step that holds $J_\nu$ fixed
while the chemistry is substepped in each cell, with the radiation field
recomputed once the ionization has advanced enough to change the opacity.
An ionization front that advances while a source brightens, or recedes
after it switches off, is a reachable target.  The hard part is the stiff
integrator and the step control, not the infrastructure.

\subsection{Cooling flows: depends on which cooling flow}

The term carries two meanings, and they price very differently.

In the \Mappings{} sense --- a parcel of gas that cools radiatively as it
moves down a prescribed density and velocity history --- a cooling flow is
a special case of time-dependent cooling on a fixed trajectory: switch off
(or freeze) the photoionizing field, and integrate the non-equilibrium
ionization and the radiative losses down a temperature ramp.  Once the
time-dependent solver of the previous subsection exists, this is a modest
addition, because the cooling terms it needs --- recombination,
free--free, and the metal-line coolants through the Tier-1 fits and the
Tier-2 $n$-level solve --- are already present.  Its one real prerequisite
is the density-dependent cooling flagged as HS7 and in
\S\ref{sec:densitycool}: a cooling flow crosses coolant critical
densities, so a $\Lambda(\Te,\nelec)$ that suppresses forbidden lines
correctly is required for it to be quantitative.

In the cluster or galaxy sense --- a self-consistent inflow driven by its
own radiative losses, with gravity and hydrodynamics --- a cooling flow is
a fluid-dynamics problem and lies outside the present scope without a
hydrodynamic solver.  The comparison with \Mappings{} concerns the first
meaning.

\subsection{Radiative shocks: hard, and off the architecture}

The \Mappings{} S5 shock solves the Rankine--Hugoniot jump conditions and
then integrates the post-shock cooling zone as a one-dimensional flow ---
mass, momentum, and energy conservation along the flow coordinate with
radiative losses --- while iterating an upstream photoionized precursor
computed from the shock's own radiation to self-consistency.  Its essence
is a one-dimensional radiation-hydrodynamic flow, not a three-dimensional
transport problem.

This is the poor fit.  \MoCHII{} has no velocity field, no fluid
equations, and no flow coordinate.  The post-shock structure is gas that
moves and compresses; the static octree grid and the Monte Carlo photon
transport represent none of that, and three-dimensional Monte Carlo
transport buys almost nothing for a one-dimensional steady shock.  The
honest route is a separate one-dimensional radiative-shock module, a
companion solver in the spirit of the 1D references \MoCHII{} already uses
for its gates: it would reuse the atomic and cooling data but not the
octree or the Monte Carlo engine, and it would carry the jump conditions,
an adaptive cooling-zone integrator, and the precursor iteration itself.
That is real work, largely decoupled from the transport code, and --- this
is the key qualification --- it would be a 1D shock, not a 3D radiative
shock.  A genuine three-dimensional radiative shock needs actual radiation
hydrodynamics, which is a different code, not an extension of this one.

\subsection{Summary and suggested order}

\begin{center}
\begin{tabular}{P{0.26\textwidth}P{0.16\textwidth}P{0.44\textwidth}}
\toprule
Capability & Difficulty & Why \\
\midrule
Time-dependent ionization & Moderate &
  Reuses all transport and rates; new work is a stiff cell ODE, a time
  loop, and a light curve.  Operator splitting maps onto the existing
  one-transport-for-each-iteration loop. \\
Cooling flow (\Mappings{} sense) & Moderate &
  A prescribed-trajectory special case of the time-dependent solver;
  needs the density-dependent $\Lambda(\Te,\nelec)$ of HS7. \\
Cooling flow (cluster sense) & Out of scope &
  Requires gravity and hydrodynamics. \\
Radiative shock (S5) & Hard &
  A separate 1D radiation-hydrodynamic module; the octree and Monte Carlo
  engine do not apply.  Not a 3D shock. \\
\bottomrule
\end{tabular}
\end{center}

\noindent A sensible order is: (1)~time-dependent ionization first, as the
most natural addition and the foundation for the rest; (2)~the
density-dependent cooling of HS7, needed for both PN accuracy and cooling
flows; (3)~the \Mappings{}-sense cooling flow as a thin layer on top of
(1) and~(2); and (4)~a decoupled 1D radiative-shock module last, and only
if a shock diagnostic is actually required.  The first three share the
same time-integration and cooling infrastructure and reinforce the PN and
hard-spectrum work of \S\ref{sec:hardplan}; the fourth stands apart.

%==============================================================================
\section{Implementation evidence consulted}

The principal local sources used for this audit were:

\begin{itemize}[leftmargin=2em]
\item \MoCHII{}: \code{docs/PLAN.md}, \code{docs/MoCHII\_physics.tex},
  \code{docs/MoCHII\_UserGuide.tex}, the transport/grid modules under
  \code{src/}, and the quantitative gates under \code{tests/}.
\item \Mthree{}: \code{README.md}, \code{lab/data/ATDAT.txt},
  \code{src/mastercode/montph7.f}, \code{montph7\_dust.f},
  \code{crosssections.f}, \code{mcteequi.f}, \code{teequi2.f},
  \code{equion.f}, \code{cool.f}, \code{multilevel.f},
  \code{grainpar.f}, \code{hgrains.f}, \code{hpahs.f},
  \code{dusttemp.f}, \code{timion.f}, and \code{shock2--shock5.f}.
\item \Mappings{}: \code{README.md}, \code{src/mappings.f},
  \code{photo6.f}, \code{photo7.f}, \code{teequi.f}, \code{teequi2.f},
  \code{timion.f}, \code{timtqui.f}, \code{cool.f}, \code{dusttemp.f},
  \code{shock5.f}, \code{neqc.f}, \code{coolc.f},
  \code{transferline.f}, \code{lab/mapStd.prefs},
  \code{lab/mapFull.prefs}, and \code{docs/DocumentionOverview.html}.
\end{itemize}

The public \Mthree{} source contains several historical or experimental
driver variants.  Statements in this memo refer to the main
\code{mastercode/montph7.f} and explicitly selected
\code{montph7\_dust.f} paths unless otherwise stated.
Statements about \Mappings{} refer to the connected public 5.2.1
\code{map52} P6/P7/S5/CC/NC/PP paths, not merely to routines present in
the directory.

%==============================================================================
\begin{thebibliography}{99}

\bibitem{jin2022}
Jin, Y., Kewley, L. J., \& Sutherland, R. S. 2022,
``Messenger Monte Carlo MAPPINGS V (M$^3$)--A Self-consistent,
Three-dimensional Photoionization Code,'' ApJ, 927, 37,
\url{https://doi.org/10.3847/1538-4357/ac48f3}.

\bibitem{jin2022fractal}
Jin, Y., Kewley, L. J., \& Sutherland, R. S. 2022,
``Theoretically Modeling Photoionized Regions with Fractal Geometry in
Three Dimensions,'' ApJL, 934, L1,
\url{https://doi.org/10.3847/2041-8213/ac80f3}.

\bibitem{sd1993}
Sutherland, R. S., \& Dopita, M. A. 1993,
``Cooling Functions for Low-density Astrophysical Plasmas,'' ApJS, 88, 253.

\bibitem{mappings521}
Sutherland, R. S., Dopita, M. A., Binette, L., et al.\ 2022,
``MAPPINGS V v5.2.1 Public Source Release,''
\url{https://mappings.anu.edu.au}.

\bibitem{lucy1999}
Lucy, L. B. 1999, ``Computing Radiative Equilibria with Monte Carlo
Techniques,'' A\&A, 344, 282.

\bibitem{verner1996}
Verner, D. A., Ferland, G. J., Korista, K. T., \& Yakovlev, D. G. 1996,
``Atomic Data for Astrophysics. II. New Analytic Fits for Photoionization
Cross Sections of Atoms and Ions,'' ApJ, 465, 487.

\end{thebibliography}

\end{document}
